\chapter{Số học}

\section{Số học đồng dư (Modular arithmetic)}
	\kactlimport{ModularArithmetic.h}
	\kactlimport{ModInverse.h}
	\kactlimport{ModPow.h}
	\kactlimport{ModLog.h}
	\kactlimport{ModSum.h}
	\kactlimport{ModMulLL.h}
	\kactlimport{ModSqrt.h}

\section{Tính nguyên tố và phân tích thừa số}
	\kactlimport{FastEratosthenes.h}
	\kactlimport{MillerRabin.h}
	\kactlimport{Factor.h}

\section{Tính chia hết}
	\kactlimport{euclid.h}
	% \kactlimport{Euclid.java}
	\kactlimport{CRT.h}

	\subsection{Đồng nhất thức Bézout}
	Với $a \neq 0$, $b \neq 0$, $d=\gcd(a,b)$ là số nguyên dương nhỏ nhất sao cho phương trình sau có nghiệm nguyên:
	$$ax+by=d$$
	Nếu $(x,y)$ là một nghiệm, thì tất cả các nghiệm có dạng:
	$$\left(x+\frac{kb}{\gcd(a,b)}, y-\frac{ka}{\gcd(a,b)}\right), \quad k\in\mathbb{Z}$$

	\kactlimport{phiFunction.h}

\section{Phân số}
	\kactlimport{ContinuedFractions.h}
	\kactlimport{FracBinarySearch.h}

\section{Bộ ba số Pythagoras}
 Các bộ ba Pythagoras nguyên thủy được sinh duy nhất bởi:
 \[ a=k\cdot (m^{2}-n^{2}),\ \,b=k\cdot (2mn),\ \,c=k\cdot (m^{2}+n^{2}), \]
 với $m > n > 0$, $k > 0$, $\gcd(m,n)=1$, và $m$ hoặc $n$ là số chẵn.

\section{Số nguyên tố}
	$p=962592769$ thỏa mãn $2^{21} \mid p-1$, rất hữu dụng cho NTT. Khi băm (hash)
	dùng $970592641$ (31-bit), $31443539979727$ (45-bit), $3006703054056749$
	(52-bit). Có 78498 số nguyên tố nhỏ hơn $1\,000\,000$.

	Căn nguyên thủy tồn tại modulo lũy thừa nguyên tố $p^a$, ngoại trừ $p = 2, a > 2$, và có tất cả $\phi(\phi(p^a))$ căn nguyên thủy.

\section{Ước lượng}
	$\sum_{d|n} d = O(n \log \log n)$.

	Số ước của $n$ nhiều nhất xấp xỉ 100 với $n < 5e4$, 500 với $n < 1e7$, 2000 với $n < 1e10$, 200\,000 với $n < 1e19$.

\section{Hàm Möbius}
\[
	\mu(n) = \begin{cases} 0 & n \textrm{ chứa ước chính phương}\\ 1 & n \textrm{ có số ước nguyên tố chẵn}\\ -1 & n \textrm{ có số ước nguyên tố lẻ}\\\end{cases}
\]
  Công thức nghịch đảo Möbius:
  \[ g(n) = \sum_{d|n} f(d) \Leftrightarrow f(n) = \sum_{d|n} \mu(d)g(n/d) \]
  Các công thức biến thể hữu ích:

  $ \sum_{d | n} \mu(d) = [ n = 1] $

  $ g(n) = \sum_{n|d} f(d) \Leftrightarrow f(n) = \sum_{n|d} \mu(d/n)g(d)$

  $ g(n) = \sum_{1 \leq m \leq n} f(\left\lfloor\frac{n}{m}\right \rfloor ) \Leftrightarrow f(n) = \sum_{1\leq m\leq n} \mu(m)g(\left\lfloor\frac{n}{m}\right\rfloor)$
